\emph{Two elementary tools.} (a) \emph{Negative bundles.} For a nondegenerate symmetric form on a vector bundle over a compact base, a maximal negative definite subbundle exists (the negative eigenspaces with respect to an auxiliary bundle metric), and any two are isomorphic: the space of maximal negative subspaces of a fixed fibre is contractible, and more concretely, if $N_1, N_2$ are maximal negative subspaces, the projection onto $N_1$ along its $h$-orthogonal complement $N_1^{\perp_h}$ (which is positive definite) restricts to an isomorphism $N_2 \to N_1$: its kernel on $N_2$ is $N_2 \cap N_1^{\perp_h}$, negative and positive definite at once, hence zero, and the dimensions agree. Write $N(h)$ for this bundle. (b) \emph{Continuity along a collar.} If $V_t$, $t \in [0,1]$, is a continuous family of subbundles of a fixed bundle over a compact base, then $V_0 \cong V_1$: with orthogonal projections $\pi_t$ for an auxiliary metric, $\pi_t|_{V_s} : V_s \to V_t$ is an isomorphism when $\|\pi_t - \pi_s\| < 1$, uniform continuity on $[0,1]$ gives finitely many steps, and composing them gives the isomorphism. Over the collar, transport along the lines $\{x\} \times [-1,1]$ for any connection identifies $T|_C$ with the pullback of $T|_\Sig$, so a subbundle of $T|_C$ is a continuous family of subbundles of $T|_\Sig$ and (b) applies to it. Finally, isomorphic bundles have isomorphic determinant lines, $\det(N \oplus \ell) \cong \det N \otimes \ell$ for a line bundle $\ell$, $\ell \otimes \ell$ is trivial for every real line bundle and a line bundle is trivial iff its $w_1$ vanishes (Fact \claimref{sheaf.imported-real-line-bundles}).

\emph{The negative bundle of $G$ near the wall.} Throughout, $(M,G)$ satisfies (T1) to (T3) (\claimref{sheaf.tameness}). Let $\tilde\shL$ be a line subbundle of $T^*M|_C$ extending $\shL$ and $Q$ a complement. Along $\Sig$, $Q \cong T^*M/\shL$ and $G|_Q$ is the induced form $\overline G$, nondegenerate of signature $(p{-}1,q)$ by Theorem \claimref{sheaf.degeneration-sheaf}(2); by compactness $G|_Q$ is nondegenerate on a thinner collar $C_0 = \Sig \times [-\varepsilon, \varepsilon]$. Then $T^*M = Q \oplus Q^{\perp_G}$ on $C_0$, the line $Q^{\perp_G}$ is isomorphic to $T^*M/Q \cong \tilde\shL$ and equals $\shL$ on $\Sig$, and $G$ restricted to it is the Schur complement $\hat\lambda$, which has the sign of the dying eigenvalue: positive on the physical side and negative on the far side (the signatures are $(p,q)$ and $(p{-}1,q{+}1)$ and $G|_Q$ has constant signature). Hence, with $N_Q := N(G|_Q)$, a rank-$q$ bundle on $C_0$ restricting to $N(\overline G)$ on $\Sig$,
\[
N(G) \cong N_Q \ \text{on the physical half of } C_0 \setminus \Sig, \qquad N(G) \cong N_Q \oplus \tilde\shL \ \text{on the far half}.
\]
Since $G$ is nondegenerate on each half of $C \setminus \Sig$, tool (b) carries these isomorphisms from $\Sig \times \{\pm\varepsilon\}$ to $\Sig \times \{\pm 1\}$. The check verifies the underlying inertia additivity $\operatorname{In}(G) = \operatorname{In}(G|_Q) + \operatorname{sign}\hat\lambda$ on fixed rational instances.

\emph{(1)} is Theorem \claimref{sheaf.conservation-degeneracy-absolute-form}: the totals $(p,q) + \operatorname{sig}K$ and $(p{-}1,q{+}1) + \operatorname{sig}K$ differ, and in clock terms the two sides sit at $d$ and $d - 2$, with $2 \neq 0$ in $\mathbb{Z}_8$. Both sides being nonempty is the only hypothesis, so the existence question of clause (a) is decided by the coarse data alone.

\emph{(2)} At a point, choose an $A$-orthonormal frame, so that $G$ has a symmetric Gram matrix $S$ and $h_D$ has Gram $D = \begin{psmallmatrix} S & \beta 1 \\ \beta 1 & -S \end{psmallmatrix}$. Then $D^2 = (S^2 + \beta^2) \oplus (S^2 + \beta^2)$. Where $\beta \neq 0$ this is positive definite, so $D$ is invertible; where $\beta = 0$ the point lies off $\Sig$, $S$ is invertible and $D = S \oplus (-S)$ is invertible. The map $J(\alpha,\alpha') = (\alpha', -\alpha)$ satisfies $J^{\mathsf T} D J = -D$ in such a frame, so the spectrum of $D$ is symmetric and the signature is $(n,n)$. Outside $C$, $h_D = G \perp (-G)$, whose second summand has signature $(q,p)$ on the physical side and $(q{+}1,p{-}1)$ on the far side.

\emph{(3), necessity.} Let $h$ be as stated and $V := N(h)$, a bundle over $M$ of rank $q + 1$ (on the physical exterior $h = G \perp H$ has signature $(p{+}1,q{+}1)$, and $M$ is connected). On $\Sig \times \{1\}$, $V \cong N(G) \oplus N(H) \cong N_Q|_{\Sig \times \{1\}} \oplus \underline{\R}$, and on $\Sig \times \{-1\}$, $V \cong N(G) \oplus N(K_-) = N(G) \cong N_Q|_{\Sig \times \{-1\}} \oplus \tilde\shL|_{\Sig \times \{-1\}}$. By tool (b) applied to $V$, $N_Q$ and $\tilde\shL$ over $C$, all identified with bundles over $\Sig$,
\[
N(\overline G) \oplus \underline{\R} \;\cong\; N(\overline G) \oplus \shL .
\]
Taking determinant lines, $\det N(\overline G) \cong \det N(\overline G) \otimes \shL$, and tensoring with $\det N(\overline G)$ (whose square is trivial) gives $\shL \cong \underline{\R}$, so $w_1(\shL) = 0$. The necessity half of the exchanged statement is the same computation with the two exteriors interchanged.

\emph{(3), sufficiency.} Let $w_1(\shL) = 0$ and take collar data $(k, u, \beta, \mu)$ (Definition \claimref{obj.sheaf.collar-data-coupling}). In the plane use the basis $a = \nO + \nInf$, $b = \nO - \nInf$, in which $H = \operatorname{diag}(-2, 2)$. For a smooth function $\varsigma$ with values in $[-1,1]$ set $K_\varsigma := \operatorname{diag}(-2\varsigma, 2)$ and
\[
h \;:=\; G \perp K_\varsigma + \kappa_a, \qquad \kappa_a(\alpha, a) = \mu\beta\,\alpha(u), \quad \kappa_a(\alpha, b) = 0 ,
\]
so $h = G \perp H$ where $\beta = 0$, $\varsigma = 1$ and $h = G \perp \operatorname{diag}(2,2)$ where $\beta = 0$, $\varsigma = -1$. Off $\Sig$ the Schur complement of the $G$-block, as in Lemma \claimref{sheaf.determinant-identity-blow-up-expansion}(1), gives
\[
\det h \;=\; \det G \cdot 2\,\bigl(-2\varsigma - \mu^2\beta^2 s\bigr), \qquad s = G^{-1}(e,e) = \frac1\lambda + O(1)
\]
(Lemma \claimref{sheaf.determinant-identity-blow-up-expansion}(2)). Let $S$ bound $|s|$ on the support of $\beta$ outside a fixed thin neighbourhood $Z$ of $\Sig$ inside $\{\beta = 1\}$ on which $s$ has the sign of the side (finite by compactness, as in Lemma \claimref{sheaf.tame-wall-induction}; $s$ may change sign outside $Z$), and choose $\mu$ with $\mu^2 S < 2$. Put $Z' := \Sig \cup (\{|s| > 2/\mu^2\} \cap Z)$, a neighbourhood of $\Sig$ inside $Z \subset \{\beta = 1\}$ (since $|s| \to \infty$ at the wall, uniformly by compactness); and take $\varsigma = 1$ on the physical side outside $Z'$, $\varsigma = -1$ on the far side outside $Z'$. Then $-2\varsigma - \mu^2\beta^2 s \neq 0$ off $\Sig$: where $\beta = 0$ it is $\mp 2$; on the support of $\beta$ outside $Z$ it is at most $-2 + \mu^2 S < 0$ (physical) or at least $2 - \mu^2 S > 0$ (far); on $Z \setminus Z'$, $\beta = 1$ and $s$ has the side's sign, giving $-2 - \mu^2 s < 0$ or $2 + \mu^2|s| > 0$; on $Z' \setminus \Sig$, $|\varsigma| \le 1$ and $\mu^2|s| > 2$ give $-2\varsigma - \mu^2 s \le 2 - \mu^2 s < 0$ (physical) and $\ge -2 + \mu^2|s| > 0$ (far). On $\Sig$, as in Lemma \claimref{sheaf.wall-gram}, $T = V_1 \oplus \operatorname{span}(k, a, b)$ orthogonally, $h|_{V_1} = \overline G$ is nondegenerate, and the block of $(k, a, b)$ has Gram $\begin{psmallmatrix} 0 & \mu & 0 \\ \mu & -2\varsigma & 0 \\ 0 & 0 & 2 \end{psmallmatrix}$ with determinant $-2\mu^2 \neq 0$ for every value of $\varsigma$. So $h$ is nondegenerate everywhere, of signature $(p{+}1,q{+}1)$, standard on the physical exterior and positive definite in the plane on the far exterior. The check runs this construction exactly along a model transversal on which $s$ changes sign inside the collar, including the interpolation zone and the wall, and confirms that it degenerates when $\mu^2 S \ge 2$.

\emph{(3), sufficiency of the exchanged statement.} $-G$ is a tame corank-one degeneration with the same radical line $\shL$, whose physical side (where its dying eigenvalue is positive) is the far side of $G$, of signature $(q, p)$ there. The construction applied to $-G$ gives $h'$ on $T$, nondegenerate, with $h'|_{T^*M} = -G$, equal to $(-G) \perp H$ outside $C$ on the far side of $G$ and to $(-G) \perp K'$ with $K'$ positive definite on the physical side. Then $h := -h'$ restricts to $G$ on $T^*M$, equals $G \perp (-H)$ on the far exterior and $G \perp (-K')$, with $-K'$ negative definite, on the physical exterior; and $-H \cong H$ by $\nInf \mapsto -\nInf$.

\emph{(4).} Sufficiency is Theorem \claimref{sheaf.one-sided-completion}. Necessity: let $h$ be as stated and $V = N(h)$ over $\Mp$, of rank $q + 1$. On $\Sig \times \{1\}$, $V \cong N_Q \oplus \underline{\R}$ as in (3). At $\Sig$, $h|_{T^*M} = G$ makes $\shL$ isotropic and $h$-orthogonal to $T^*M$; with $V_1$ a complement of $\shL$ in $T^*M|_\Sig$, $h|_{V_1} = \overline G$ is nondegenerate and $T|_\Sig = V_1 \oplus W$, $W := V_1^{\perp_h}$ of rank three, nondegenerate, containing $\shL$, of signature $(p{+}1,q{+}1) - (p{-}1,q) = (2,1)$. In $W$, $\shL^{\perp}$ has rank two and radical $\shL$, so $\ell := \shL^\perp/\shL$ is a nondegenerate line and $W = \ell' \perp P$ with $\ell' \cong \ell$ and $P := \ell'^{\perp}$ a plane of signature $(1,1)$ containing the isotropic line $\shL$; the signature count forces $\ell$ positive. The negative line $N(W) = N(P)$ meets $\shL$ trivially, so it maps isomorphically onto $P/\shL$, which $h$ pairs nondegenerately with $\shL$ ($\shL^{\perp_P} = \shL$): $N(W) \cong P/\shL \cong \shL^* \cong \shL$, the partner line of the residue theorem (\claimref{sheaf.metrics-residue-square-root-twist}) in this setting. Hence $V|_\Sig \cong N(\overline G) \oplus \shL$. Tool (b) on the physical half of $C$ gives $N(\overline G) \oplus \shL \cong N(\overline G) \oplus \underline{\R}$, and the determinant argument of (3) gives $w_1(\shL) = 0$. The check verifies on the wall Gram of Lemma \claimref{sheaf.wall-gram} that the isotropic line's orthogonal is positive, so the negative line pairs with $\shL$.

\emph{(5).} The same argument with $\Lambda \oplus \Lambda^*$ in place of $\underline{\R}^2$: the negative line of the hyperbolic plane $\Lambda \oplus \Lambda^*$ meets the isotropic line $\Lambda$ trivially and so is isomorphic to $(\Lambda \oplus \Lambda^*)/\Lambda = \Lambda^* \cong \Lambda$; hence $V|_{\Sig \times \{1\}} \cong N_Q \oplus \Lambda$, and the comparison at $\Sig$ gives $\det N(\overline G) \otimes \shL \cong \det N(\overline G) \otimes \Lambda|_\Sig$, that is $\shL \cong \Lambda|_\Sig$. A line bundle on $\Sig$ extends over $\Mp$ exactly when its class lies in the image of the restriction on $H^1(\,\cdot\,; \mathbb{Z}_2)$, line bundles being classified by $w_1$ (Fact \claimref{sheaf.imported-real-line-bundles}).

\emph{The consequence.} By (1), clause (a)'s condition fails on every wall satisfying the hypotheses (both sides nonempty), whatever the torsion; by (2), dropping the single exterior model leaves nothing to obstruct; by (3) and (4), at one rung the obstruction is exactly $w_1(\shL)$. Along the wall sequence of the untwisted completion the pair $W(\Sig)$ is carried unchanged: the rung shifts $(p,q)$ by $(1,1)$ and preserves both clock positions, and the shadow wall's radical line has the global spanning section of the proof of the infinite wall sequence, so its $w_1$ is $0$ (Theorem \claimref{sheaf.infinite-wall-sequence}), as is that of $\shL$ under the standing hypothesis; conservation for twisted sequences is open.
